% !PRINTED 9
\weberpage{9}
%<<CONT
ments of a continuous ordered set $\mathfrak M$ in which addition
exists, and if $c>a$, then one obtains a cut $(A,B)$ in $\mathfrak M$
by assigning to $A$ all elements $x$ for which $a+x\le c$, and to $B$
those for which $a+x>c$. This cut is produced by an element $b$, for
which $a+b=c$ must hold.

The natural numbers form, according to our definition, a measurable set
in which a least element, namely $1$, occurs. The manifold of rational
fractions is likewise measurable when addition and subtraction are
explained according to the known rules of fractional calculation.
Especially important, and as it were typical for the measurable sets,
is the manifold of rectilinear segments, or lengths of a line, which
are added simply by laying them end to end. Quantities of material,
compared by means of the balance, and intervals of time, measured by
the clock, also furnish examples of measurable sets. The manner of
measuring does not lie in the nature of the manifold itself, but is put
into it by the thinking observer; thus, for example, it would be just as
admissible to understand by the sum $a+b$ of two segments $a$ and $b$
the hypotenuse of a right triangle with legs $a,b$, instead of, as is
usually assumed, the segment composed from $a$ and $b$ by laying them
end to end.

In order to make the continuous set $\mathfrak S$ of cuts in the
manifold $\mathfrak R$ of rational fractions into a measurable one,
observe first the following.

If $\alpha=(A,B)$ is an element in $\mathfrak S$ and $\mu$ is an
arbitrarily given rational fraction, then one can always choose an
element $a'$ in $A$ so that $a'+\mu=b'$ is contained in $B$. For if one
chooses two arbitrary elements $a,b$, then one can determine the
natural number $m$ so that $m\mu>b-a$, hence $a+m\mu$ is contained in
$B$. If then $h$ is the smallest integer for which $a+h\mu$ is
contained in $B$, then $a+(h-1)\mu=a'$ is in $A$ and hence
$a'+\mu=b'$ is in $B$.

We now understand by the sum $(A,B)+(A',B')=(A'',B'')$, or
$\alpha+\alpha'=\alpha''$, the cut in $\mathfrak R$ which is obtained
by assigning a rational fraction $a''$ to $A''$ only when an $a$ in
$A$ and an $a'$ in $A'$ exist such that $a''\le a+a'$. In fact
$(A'',B'')$ is a cut; for if $a''$ is contained in $A''$, the same
holds for every smaller fraction, and there are fractions which are in
$A''$ and fractions which are not in
%CONT>>
